\changecolours{ScoutGreen}
\section{Conclusão \& Próximos Passos}

% ------------------------------------------------------------------------------
\twocolframe[title={Síntese da Aula \& Orientações Avaliativas},leftcol=7cm,rightcol=7cm]{%
  \alert{O Que Aprendemos Hoje}\par
  \itemseps[topsepi=1mm,itemsepi=1.5mm,labelsep=2mm,leftmargini=3.5mm]
  \begin{itemize}
    \item \textbf{Ciclo de Vida do Projeto:} As quatro fases indispensáveis para a engenharia de dados (Requisitos, Conceitual, Lógico e Físico).
    \item \textbf{O MER de Peter Chen (1976):} Concepção ontológica de entidades e conjuntos de entidades.
    \item \textbf{Classificação Completa de Atributos:} Identificadores, simples, compostos, multivalorados, armazenados e derivados.
    \item \textbf{A Ponte Conceitual-Relacional:} Como atributos complexos do MER são adaptados para respeitar a 1FN de Codd.
  \end{itemize}
}{%
  \alert{Atividades e Próximos Passos}\par
  \itemseps[topsepi=1mm,itemsepi=1.5mm,labelsep=2mm,leftmargini=3.5mm]
  \begin{itemize}
    \item \textbf{Lista de Exercícios 01 (N1):} O Bloco 3 está disponível no plano pedagógico e encerra o conteúdo cobrado na primeira lista avaliativa.
    \item \textbf{Aula 04 (Sábado Letivo - 10/10/2026):} Encontro assíncrono dedicado à resolução e entrega final da Lista 01 (N1).
    \item \textbf{Prévia da Aula 05:} Relacionamentos no MER, grau, papéis, mapeamento de cardinalidades (1:1, 1:N, N:M) e restrições de participação total e parcial.
  \end{itemize}
}
